% part: intuitionistic-logic
% chapter: introduction
% section: natural-deduction

\documentclass[../../../include/open-logic-section]{subfiles}

\begin{document}

\olfileid{int}{int}{ntd}

\olsection{સ્વાભાવિક નિગમન}

$\FalseCl$ નિયમો વિનાનું સ્વાભાવિક નિગમન સ્ફુરણાવાદી તર્ક
માટેનું પ્રમાણભૂત !!{derivation} તંત્ર છે. અહીં નિયમો ફરી
આપીએ છીએ અને BHK અર્થઘટન વડે તેમની પ્રેરણા સમજાવીએ છીએ.
દરેક કિસ્સામાં નિયમ એવો છે કે આધારવાક્યોની રચનાઓ હોય, તો
નિષ્કર્ષની પણ રચના હોય એવું કહી શકાય.

સ્વાભાવિક નિગમનની !!{derivation}sમાં નિવૃત્ત ન કરેલી ધારણાઓ
હોતી હોવાથી, એવી !!a{derivation}ને---દાખલા તરીકે
!!{undischarged} ધારણાઓ~$\Gamma$માંથી $!A$ની નિષ્પત્તિને---
એવા વિધેય તરીકે જોવી જોઈએ જે દરેક $!B \in \Gamma$ની
રચનાઓને~$!A$ની રચનામાં ફેરવે છે. !!{undischarged} ધારણાઓ
વિના~$!A$ની !!a{derivation} હોય, તો BHK અર્થઘટનના અર્થમાં
~$!A$ની રચના હોય છે. જોકે ચર્ચામાં જરૂર ન હોય ત્યારે
~$\Gamma$નો ઉલ્લેખ છોડીશું.

ધારણા~$!A$ પોતે જ !!{undischarged} ધારણા~$!A$માંથી~$!A$ની
!!a{derivation} છે. આ BHK અર્થઘટનને અનુરૂપ છે: રચનાઓ પરનું
તાદાત્મ્ય વિધેય~$!A$ની કોઈ પણ રચનાને~$!A$ની રચનામાં
ફેરવે છે.

\subsection{સંયોજન}

\begin{defish}
\AxiomC{$!A$}
\AxiomC{$!B$}
\RightLabel{\Intro{\land}}
\BinaryInfC{$!A \land !B$}
\DisplayProof
\hfill
\begin{tabular}{l}
\AxiomC{$!A \land !B$}
\RightLabel{\Elim{\land}}
\UnaryInfC{$!A$}
\DisplayProof
\\[3ex]
\AxiomC{$!A \land !B$}
\RightLabel{\Elim{\land}}
\UnaryInfC{$!B$}
\DisplayProof
\end{tabular}
\end{defish}

ધારો કે $!A_1$ અને $!A_2$ની રચનાઓ અનુક્રમે $N_1$, $N_2$
છે. તો $!A_1 \land !A_2$ની રચના, એટલે જોડું
$\tuple{N_1, N_2}$, પણ આપણી પાસે છે.

BHK અર્થઘટન મુજબ $!A_1 \land !A_2$ની રચના એ જોડું
$\tuple{N_1, N_2}$ છે. \sourcecorrection{OLMD-085}{મૂળમાં
અહીં $!A_1 \land !A_1$ છે, જ્યારે આગળ $N_1$ અને $N_2$
અનુક્રમે $!A_1$ તથા $!A_2$ની રચનાઓ છે. આ સંયોજન
$!A_1 \land !A_2$ કર્યું છે.} આવું જોડું હોય, તો દરેક
સંયોજ્યની રચના પણ હોય છે: $N_1$ એ~$!A_1$ની અને $N_2$
એ $!A_2$ની રચના છે.

\subsection{શરતી સંયોજક}

\begin{defish}
  \AxiomC{$\Discharge{!A}{u}$}
  \DeduceC{$!B$}
  \DischargeRule{$\Intro{\lif}$}{u}
  \UnaryInfC{$!A \lif !B$}
  \DisplayProof
\hfill
  \AxiomC{$!A \lif !B$}
  \AxiomC{$!A$}
  \RightLabel{$\Elim{\lif}$}
  \BinaryInfC{$!B$}
  \DisplayProof
\end{defish}

!!{undischarged} ધારણા~$!A$માંથી~$!B$ની !!a{derivation}
હોય, તો એવું વિધેય~$f$ હોય છે જે~$!A$ની રચનાઓને~$!B$ની
રચનાઓમાં ફેરવે છે. આ જ વિધેય $!A \lif !B$ની રચના છે.
તેથી $\Intro\lif$ના આધારવાક્યની રચના~$!A$ની રચના પર
આધારિત હોય, તો નિષ્કર્ષ $!A \lif !B$ની રચના હોય છે.

બીજી તરફ, ધારો કે $!A$ની રચના~$N$ અને $!A \lif !B$ની
રચના~$f$ હોય. $!A \lif !B$ની રચના એવું વિધેય છે જે
$!A$ની રચનાઓને~$!B$ની રચનાઓમાં ફેરવે છે. તેથી $f(N)$
એ~$!B$ની રચના છે; એટલે $\Elim\lif$ના નિષ્કર્ષની રચના છે.

\subsection{વિયોજન}

\begin{defish}
\begin{tabular}{l}
\AxiomC{$!A$}
\RightLabel{\Intro{\lor}}
\UnaryInfC{$!A \lor !B$}
\DisplayProof
\\[3ex]
\AxiomC{$!B$}
\RightLabel{\Intro{\lor}}
\UnaryInfC{$!A \lor !B$}
\DisplayProof
\end{tabular}
\hfill
\AxiomC{$!A \lor !B$}
\AxiomC{$\Discharge{!A}{n}$}
\DeduceC{$!C$}
\AxiomC{$\Discharge{!B}{n}$}
\DeduceC{$!C$}
\DischargeRule{\Elim{\lor}}{n}
\TrinaryInfC{$!C$}
\DisplayProof
\end{defish}

~$!A_i$ની રચના~$N_i$ હોય, તો તેમાંથી~$!A_1 \lor !A_2$ની
રચના~$\tuple{i, N_i}$ બનાવી શકીએ. બીજી તરફ, ધારો કે
~$!A_1 \lor !A_2$ની રચના છે, એટલે જોડું $\tuple{i, N_i}$
છે જેમાં $N_i$ એ~$!A_i$ની રચના છે. વળી વિધેયો $f_1$,
$f_2$ અનુક્રમે $!A_1$, $!A_2$ની રચનાઓને~$!C$ની
રચનાઓમાં ફેરવે છે. તો $f_i(N_i)$ એ~$!C$ની રચના છે,
જે~$\Elim\lor$નો નિષ્કર્ષ છે.

\subsection{અસંગતિ}

\begin{defish}
  \AxiomC{$\lfalse$}
  \RightLabel{\FalseInt}
  \UnaryInfC{$!A$}
  \DisplayProof
\end{defish}

!!{undischarged} ધારણાઓ~$!B_1$, \dots, $!B_n$માંથી
~$\lfalse$ની !!a{derivation} હોય, તો વિધેય~$f(M_1,
\dots, M_n)$ એ $!B_1$, \dots, $!B_n$ની રચનાઓને
~$\lfalse$ની રચનામાં ફેરવે છે. $\lfalse$ની કોઈ રચના
ન હોવાથી, $!B_1$, \dots, $!B_n$ બધાની રચનાઓ પણ એકસાથે
હોઈ શકતી નથી. તેથી $f$માં એ ગુણધર્મ પણ છે કે \emph{જો}
$M_1$, \dots, $M_n$ અનુક્રમે $!B_1$, \dots, $!B_n$ની
રચનાઓ હોય, \emph{તો} $f(M_1, \dots, M_n)$ એ~$!A$ની
રચના છે.

\subsection{$\lnot$ માટેના નિયમો}

$\lnot !A$ને $!A \lif \lfalse$ તરીકે વ્યાખ્યાયિત કર્યું
હોવાથી કડક અર્થમાં~$\lnot$ માટે અલગ નિયમો જરૂરી નથી.
તેમ છતાં આપીએ તો તે નીચે જેવા હશે:

\begin{defish}
\AxiomC{$\Discharge{!A}{n}$}
\noLine
\DeduceC{$\lfalse$}
\DischargeRule{\Intro{\lnot}}{n}
\UnaryInfC{$\lnot !A$}
\DisplayProof
\hfill
\AxiomC{$\lnot !A$}
\AxiomC{$!A$}
\RightLabel{\Elim{\lnot}}
\BinaryInfC{$\lfalse$}
\DisplayProof
\end{defish}


\subsection{\usetoken{P}{derivation}નાં ઉદાહરણો}

\begin{enumerate}
\item $\Proves !A \lif (\lnot !A \lif \lfalse)$, એટલે કે $\Proves !A
  \lif ((!A \lif \lfalse) \lif \lfalse)$
  \begin{prooftree}
    \AxiomC{$\Discharge{!A}{2}$}
    \AxiomC{$\Discharge{!A \lif \lfalse}{1}$}
    \RightLabel{\Elim\lif}
    \BinaryInfC{$\lfalse$}
    \DischargeRule{\Intro\lif}{1}
    \UnaryInfC{$(!A \lif \lfalse)\lif \lfalse$}
    \DischargeRule{\Intro\lif}{2}
    \UnaryInfC{$!A \lif (!A \lif \lfalse) \lif \lfalse$}
  \end{prooftree}

\item $\Proves ((!A \land !B) \lif !C) \lif (!A \lif (!B \lif !C))$
  \begin{prooftree}
    \AxiomC{$\Discharge{(!A \land !B) \lif !C}{3}$}
    \AxiomC{$\Discharge{!A}{2}$}
    \AxiomC{$\Discharge{!B}{1}$}
    \RightLabel{\Intro\land}
    \BinaryInfC{$!A \land !B$}
    \RightLabel{\Elim\lif}
    \BinaryInfC{$!C$}
    \DischargeRule{\Intro\lif}{1}
    \UnaryInfC{$!B \lif !C$}
    \DischargeRule{\Intro\lif}{2}
    \UnaryInfC{$!A \lif (!B \lif !C)$}
    \DischargeRule{\Intro\lif}{3}
    \UnaryInfC{$((!A \land !B) \lif !C) \lif (!A \lif (!B \lif !C))$}
  \end{prooftree}

\item $\Proves \lnot(!A \land \lnot !A)$, એટલે કે $\Proves (!A \land (!A
  \lif \lfalse))\lif \lfalse$
  \begin{prooftree}
    \AxiomC{$\Discharge{!A \land (!A \lif \lfalse)}{1}$}
    \RightLabel{\Elim\land}
    \UnaryInfC{$!A \lif \lfalse$}
    \AxiomC{$\Discharge{!A \land (!A \lif \lfalse)}{1}$}
    \RightLabel{\Elim\land}
    \UnaryInfC{$!A$}
    \RightLabel{\Elim\lif}
    \BinaryInfC{$\lfalse$}
    \DischargeRule{\Intro\lif}{1}
    \UnaryInfC{$(!A \land (!A \lif \lfalse))\lif \lfalse$}
  \end{prooftree}

\item $\Proves \lnot\lnot(!A \lor \lnot !A)$, એટલે કે $\Proves ((!A \lor
  (!A \lif \lfalse))\lif \lfalse)\lif \lfalse$
  \begin{prooftree}\hskip-3em
    \AxiomC{$\Discharge{(!A \lor (!A \lif \lfalse))\lif \lfalse}{2}$}
    \AxiomC{$\Discharge{(!A \lor (!A \lif \lfalse))\lif \lfalse}{2}$}
    \AxiomC{$\Discharge{!A}{1}$}
    \RightLabel{\Intro\lor}
    \UnaryInfC{$!A \lor (!A \lif \lfalse)$}
    \RightLabel{\Elim\lif}
    \BinaryInfC{$\lfalse$}
    \DischargeRule{\Intro\lif}{1}
    \UnaryInfC{$!A \lif \lfalse$}
    \RightLabel{\Intro\lor}
    \UnaryInfC{$!A \lor (!A \lif \lfalse)$}
    \RightLabel{\Elim\lif}
    \insertBetweenHyps{\hskip -4em}
    \BinaryInfC{$\lfalse$}
    \DischargeRule{\Intro\lif}{2}
    \UnaryInfC{$((!A \lor (!A \lif \lfalse))\lif \lfalse)\lif \lfalse$}
  \end{prooftree}
\end{enumerate}

\begin{prop}
  સ્ફુરણાવાદી તર્કના સ્વાભાવિક નિગમનમાં $\Gamma \Proves !A$
  હોય, તો પ્રશિષ્ટ તર્કમાં પણ $\Gamma \Proves !A$. ખાસ કરીને,
  $!A$ સ્ફુરણાવાદી પ્રમેય હોય તો પ્રશિષ્ટ પ્રમેય પણ છે.
\end{prop}

\begin{proof}
  સ્વાભાવિક નિગમનનો દરેક નિયમ પ્રશિષ્ટ સ્વાભાવિક નિગમનમાં
  પણ નિયમ છે. તેથી સ્ફુરણાવાદી તર્કની દરેક !!{derivation}
  પ્રશિષ્ટ તર્કમાં પણ !!a{derivation} છે.
\end{proof}

\begin{prob}
  નીચેનાં !!{formula}sની સ્ફુરણાવાદી તર્કમાં
  !!{derivation}s આપો:
  \begin{enumerate}
    \item $(\lnot !A \lor !B) \lif (!A \lif !B)$
    \item $\lnot\lnot\lnot !A \lif \lnot !A$
    \item $\lnot\lnot (!A \land !B) \liff (\lnot\lnot !A\land \lnot \lnot !B)$
    \item $\lnot (!A \lor !B) \liff (\lnot !A \land \lnot !B)$
    \item $(\lnot !A \lor \lnot !B) \lif \lnot (!A \land !B)$
    \item $\lnot\lnot ( !A \land !B) \lif  (\lnot\lnot !A \lor
    \lnot\lnot !B)$
  \end{enumerate}
\end{prob}

\end{document}
